% Resultado exploratorio: selecao e avaliacao nos mesmos dez vasos.
\subsection{Busca combinatória dos sensores MQ com validação LOSO}

Para avaliar uma configuração sem os três sensores ambientais, foram testados
todos os 63 subconjuntos não vazios dos seis MQ brutos. Cada subconjunto foi
avaliado por validação \textit{leave-one-subject-out} (LOSO) dos dez vasos
biológicos. Em cada rodada, o modelo Extra Trees, com 900 árvores e
hiperparâmetros fixos, foi treinado nos outros nove vasos e testado em todas
as coletas do vaso excluído. As probabilidades das leituras foram agregadas
pela média para produzir uma única previsão por coleta. A métrica de seleção
foi a acurácia convencional nas 36 coletas reunidas das dez rodadas.

\begin{table}[htbp]
\centering
\caption{Combinações MQ selecionadas da busca exaustiva por LOSO, avaliadas
por acurácia convencional por coleta.}
\label{tab:busca_mq_loso}
\begin{tabular}{lrr}
\hline
Sensores MQ & Coletas corretas & Acurácia (\%) \\
\hline
MQ2, MQ3, MQ135, MQ138 & 30/36 & 83,33 \\
MQ2, MQ3 & 29/36 & 80,56 \\
MQ2, MQ3, MQ8 & 29/36 & 80,56 \\
MQ2, MQ3, MQ8, MQ135, MQ138 & 29/36 & 80,56 \\
Todos os seis MQ & 28/36 & 77,78 \\
\hline
\end{tabular}
\end{table}

MQ2, MQ3, MQ135 e MQ138 formaram o único subconjunto com 30 acertos em 36
coletas, sendo 15 de 17 coletas inoculadas e 15 de 19 coletas saudáveis.
O candidato anterior que incluía MQ2, MQ135, MQ138 e as três variáveis
ambientais havia alcançado 32/36 coletas no mesmo protocolo. A busca dos MQ
reduziu a diferença para duas coletas, sem usar entradas ambientais. Como o
subconjunto vencedor foi escolhido após comparar 63 alternativas nos mesmos
dez vasos, esta estimativa ainda pode favorecer a configuração selecionada;
novos vasos são necessários para confirmar seu desempenho.
